\par\Needspace{8\baselineskip}\bigskip
\SourceLeaf{457}
\pdfbookmark[3]{Exposition, 9}{pallas-II-article9}
\begin{center}9.\end{center}
\noindent\hspace*{\parindent}Supposons que les perturbations ajoutent aux équations fondamentales les quantités \(\xi,\eta,\zeta\), de sorte que nous ayons
\[
\begin{aligned}
0&=\dfrac{\mathrm{dd}x}{\mathrm{d}t^2}+\dfrac{mx}{r^3}+\xi\\[5pt]
0&=\dfrac{\mathrm{dd}y}{\mathrm{d}t^2}+\dfrac{my}{r^3}+\eta\\[5pt]
0&=\dfrac{\mathrm{dd}z}{\mathrm{d}t^2}+\dfrac{mz}{r^3}+\zeta.
\end{aligned}
\]
\par On voit, que \(\xi,\eta,\zeta\) sont les forces résultantes de la décomposition de toutes celles, qui troublent le mouvement relatif de la planète par rapport au soleil, suivant des directions opposées à celles des trois axes. On peut aussi imaginer ces forces réunies en une seule, dont l’intensité soit \(=\chi\) et la direction parallèle à un rayon de notre sphère mené du centre vers le point de la surface \(Q\). On aura donc
\[
\begin{aligned}
\xi&=-\chi\cos XQ\\
\eta&=-\chi\cos YQ\\
\zeta&=-\chi\cos ZQ.
\end{aligned}
\]
\par Il peut être utile d’introduire au lieu des forces \(\xi,\eta,\zeta\) celles qui naissent de la décomposition de \(\chi\) suivant trois autres directions perpendiculaires entre elles. Supposons d’abord que la première agisse perpendiculairement au plan de l’orbite osculatrice, la seconde dans la direction du rayon vecteur, et par conséquent la troisième perpendiculairement au rayon vecteur dans le plan de l’orbite, ou plus exactement, que ces trois résultantes agissent parallèlement aux trois rayons de la sphère menés vers les points \(M,L,\Lambda\), de sorte qu’elles deviennent
\[
\begin{aligned}
\chi\cos MQ&=-\xi\cos XM-\eta\cos YM-\zeta\cos ZM\\
\chi\cos LQ&=-\xi\cos XL-\eta\cos YL-\zeta\cos ZL\\
\chi\cos\Lambda Q&=-\xi\cos X\Lambda-\eta\cos Y\Lambda-\zeta\cos Z\Lambda.
\end{aligned}
\]
\par Ainsi la première sera
\[
=-\dfrac{A\xi+B\eta+C\zeta}{\BareRoot{mp}}\mathpunct{\raisebox{.45ex}{,}}
\]
\SourceLeaf{458}la seconde
\[
=-\dfrac{x\xi+y\eta+z\zeta}{r}\mathpunct{\raisebox{.45ex}{,}}
\]
la troisième
\[
=-\dfrac{(Bz-Cy)\xi+(Cx-Az)\eta+(Ay-Bx)\zeta}{r\BareRoot{mp}}\mathpunct{\cdot}
\]
Nous désignerons ces trois forces par \(mW,mT,mV\).
